\editionPage[\pdfbookmark[0]{Fiches outils}{bm-outils-debut-parent}]{outils-debut}{Fiche outil}{Conduire un calcul}{À consulter au fil des chapitres.}
\begin{blocEdition}{outils-debut}
\editionEtape{Choisir}Identifier la question, les données et leurs unités. Faire une estimation avant de lancer le calcul.\editionEtape{Organiser}Parenthèses, puissances 2 ou 3, produits et quotients, sommes et différences. À priorité égale, garder l’ordre gauche-droite.\editionEtape{Écrire}Une ligne par transformation ; conserver les parties inchangées. Ne pas employer le signe égal entre deux quantités différentes.\editionEtape{Contrôler}Addition inverse, produit inverse, ordre de grandeur, somme des parts, signe et unité. Conserver la valeur exacte jusqu’à l’arrondi final.\editionRetenir{Un résultat ne termine le problème qu’après son interprétation. Pour des boîtes ou des places, le contexte impose parfois un nombre entier.}
\end{blocEdition}
\editionPage{outils-2}{Fiche outil}{Construire et justifier en géométrie}{À consulter au fil des chapitres.}
\begin{blocEdition}{outils-2}
\editionEtape{Préparer}Faire une figure de recherche, écrire les données et les coder. Une figure demandée se refait sur feuille avec les instruments utiles.\editionEtape{Choisir}Règle pour l’alignement, compas pour le report de longueurs, équerre pour les perpendiculaires, rapporteur pour les angles.\editionEtape{Raisonner}Distinguer les données de la conclusion. Une mesure suggère une conjecture ; une propriété avec ses hypothèses permet de prouver.\editionEtape{Relire}Vérifier le sommet d’un angle, les points du milieu, la base relative à une hauteur et l’ordre des sommets d’un quadrilatère.\editionRetenir{Hauteur : sommet et perpendicularité. Médiane : sommet et milieu. Médiatrice : milieu et perpendicularité.}
\end{blocEdition}
\editionPage{outils-3}{Fiche outil}{Lire un graphique ou un tableau}{À consulter au fil des chapitres.}
\begin{blocEdition}{outils-3}
\editionEtape{Identifier}Repérer les grandeurs, les unités, les axes, l’origine et les pas de graduation. Vérifier si un axe est tronqué.\editionEtape{Lire}Une abscisse puis une ordonnée forment un couple ; ne pas les permuter. Un point du tableau correspond à une entrée et à sa sortie.\editionEtape{Interpréter}Distinguer effectif et fréquence, moyenne et valeurs observées. Ne pas prolonger une tendance au-delà du domaine sans justification.\editionEtape{Proportionnalité}Un coefficient identique dans tout le tableau, ou des points alignés avec l’origine, caractérise la proportionnalité.\editionRetenir{Le choix des axes et de la représentation peut aider ou tromper. Vérifier ce que le graphique montre réellement.}
\end{blocEdition}
\editionPage{outils-4}{Fiche outil}{Tableur et blocs : explorer avec méthode}{À consulter au fil des chapitres.}
\begin{blocEdition}{outils-4}
\editionEtape{Tableur}Une formule commence par \texttt{=}. Pour $3x+2$ avec $x$ en A2 : \texttt{=3*A2+2}. Une moyenne de données en A2 à A6 : \texttt{=MOYENNE(A2:A6)}.\editionEtape{Expérimenter}Tester plusieurs entrées, dont zéro si le contexte le permet. Une liste de tests peut faire naître une conjecture ; chercher ensuite une justification.\editionEtape{Programmer}Repérer entrée, expression et sortie. Pour un tracé, préciser la direction initiale et le stylo. Dans une boucle, repérer exactement les instructions répétées.\editionEtape{Contrôler}Prévoir une sortie à la main puis comparer. Pour un polygone, compter côtés, longueurs et changements de direction.\editionRetenir{En Cinquième, on lit les données saisies et on travaille des répétitions simples à nombre fixé. Les conditions et les mises à jour de variables ne sont pas nécessaires à ces activités.}
\end{blocEdition}
